%% =====================================================================
%%  T-23-09  The Fifty-Four Percent Ceiling
%%  -------------------------------------------------------------------
%%  One idea per file. Core is mandatory. Slots in canonical order.
%%  Max 700 words. No layout commands. No filler. New source material
%%  for this entry goes in sources/B8/T-23-09.tex -- never by
%%  rewriting this file (docs/RULES.md R0.1).
%% =====================================================================
%% @kind: topic
%% @id: T-23-09
%% @title: The Fifty-Four Percent Ceiling
%% @subtitle: Unaided lie detection barely beats chance --- and the errors are not symmetric
%% @chapter: c23
%% @order: 90
%% @status: stable
%% @sources: B8
%% @updated: 2026-09-20

\topic{T-23-09}{The Fifty-Four Percent Ceiling}{Unaided lie detection barely beats chance --- and the errors are not symmetric}

\begin{core}
B8's meta-analysis measures the whole detection enterprise: across 206 studies and 24,483 judges, unaided real-time lie-truth judgment lands at 53.98 percent correct. The errors are lopsided --- 61.34 percent of truths are correctly accepted, but only 47.55 percent of lies are caught. The average reader is not a lie detector; they are a truth accepter with a small edge.
\end{core}
\begin{mechanism}
The truth bias is the default setting: 55.23 percent of all judgments call the message true. The stereotype that is supposed to correct it --- the liar as tormented, fidgeting, gaze-averse --- is a moral caricature, and calm, rationalized deceit does not match it, so most lies pass looking ordinary. Three measured facts sharpen the picture. First, the eye is the weak channel: judges watching silent video discriminate at near chance ($d = .077$) while judges who can hear do far better (audio $d = .419$, audio plus video $d = .438$); face-only and body-only views are worthless ($d = .01$ and $-.15$). What the eye reads as guilt is mostly noise. Second, pressure backfires into the stereotype: senders motivated to be believed were judged truthful by 46.84 percent of video-only receivers, against 54.44 percent for unmotivated senders --- whether they were lying or not. The motivated truth-teller is the reading's chief casualty. Third, expertise does not compound: across 19 studies in which professionals and novices judged the same messages, the difference is nil ($d = -.025$); experts were slightly more skeptical (52.28 versus 55.69 percent truth calls) without being more right. And when real lies surface, they surface late --- over days to months, through physical evidence and third parties, not from behaviour at the moment.
\end{mechanism}
\begin{conditions}
The 54 percent figure is for unaided, real-time judgment with no instruments --- exactly a live conversation. It improves with a baseline: judges who saw the sender's ordinary manner first reached 55.91 percent against 52.26 ($d = .239$), and unrehearsed senders are easier to read than prepared ones ($d = -.144$). The ceiling is not a licence to disbelieve everything: an effect of $d \approx .40$ sits at the 60th percentile of meta-analysed social-psychological findings --- real, but far too thin to carry a verdict alone.
\end{conditions}
\begin{application}
  \item Price every behavioural read at 54 percent, and act on verification rather than manner (T-23-03).
  \item Privilege the ear over the eye; treat silent-channel impressions --- photographs, silent video, a frozen face --- as barely better than guessing.
  \item Sample a baseline before reading anyone: exposure to the ordinary is the meta-analysis's one receiver-side lever (T-23-05).
  \item Discount guilt theatre: high stakes make sincerity itself look staged. Read questions, clusters and timing --- not discomfort.
  \item Expect discovery to run on records and witnesses, not live reads; log what is agreed (T-24-06).
\end{application}
\begin{feedback}
  \item Working: \emph{he seemed shady} registers as a 54 percent claim that still needs evidence to act on.
  \item Working: you test the read against a baseline before trusting it.
  \item Failing: you confront on manner alone.
  \item Failing: you read a nervous truth-teller as a liar because the stakes make them sweat.
\end{feedback}
\begin{failure}
The ceiling cuts both ways. A skilled liar shops for the setting: where nearly everyone tells the truth, base rates alone push detection toward 60 percent, and that is precisely where deceits belong. And the reader who promotes themselves past the ceiling --- trusting conviction over calibration --- converts a small honest edge into confident error at scale. The meta-analysis's flat expertise line is the receipt: seasoned professionals outperformed novices by nothing.
\end{failure}
\begin{countermeasures}
  \item When someone plays face-reader on you, the numbers are armour: their unaided read of you is, on the evidence, a coin with an edge. Ask what evidence stands behind the verdict.
  \item Anxiety about being believed is itself the tell they are misreading. Say so, plainly, when falsely read.
  \item Keep records, so that your honesty is verifiable after the fact --- the channel on which truth actually wins.
  \item In honest settings, weight the base rate over the read.
\end{countermeasures}

\sources{B8}
\seesources
\seealso{T-23-04, T-23-05, T-23-03}
